\documentclass[11pt]{article}

\usepackage[margin=1in]{geometry}
\usepackage{amsmath,amssymb,amsthm}
\usepackage{mathtools}
\usepackage{enumitem}
\usepackage[colorlinks=true,linkcolor=blue,citecolor=blue,urlcolor=blue]{hyperref}

\newtheorem{theorem}{Theorem}
\newtheorem{lemma}{Lemma}
\newtheorem{corollary}{Corollary}
\theoremstyle{definition}
\newtheorem{definition}{Definition}
\newtheorem{remark}{Remark}

\newcommand{\Z}{\mathbb{Z}}
\newcommand{\N}{\mathbb{N}}
\newcommand{\nmin}{n_{\min}}
\DeclareMathOperator{\lcm}{lcm}

\title{\textbf{How Large a $3x+1$ Cycle Must Be}\\[2pt]
\large A self-contained derivation of cycle lower bounds via the dropping dictionary}
\author{Darcy Thomas \\ \small with computational assistance}
\date{\today}

\begin{document}
\maketitle

\begin{abstract}
We give a short, self-contained, and elementary derivation of explicit lower
bounds on any nontrivial cycle of the $3x+1$ (Collatz) map, organised around the
\emph{dropping dictionary}: the observation that one round of the map is an exact
affine substitution and that a cycle is therefore a \emph{closed word} in a finite
alphabet. From the closed-word equation we recover, by a one-line
geometric--mean inequality, the classical bound
$\nmin \le 1/(2^{K/m}-3)$ on the least element of a cycle. Combined with the
standard computational verification that the map has no cycle with least element
below $B = 2^{68}$, this forces the step ratio $K/m$ to approximate $\log_2 3$
to within $2^{-69}$, and the theory of continued fractions then forces any
nontrivial cycle to contain at least
\[
  m \ \ge\ 72{,}057{,}431{,}991 \quad\text{odd elements,}
\]
a full period of at least $1.86\times10^{11}$ steps, and a least element of at
least $2^{68}$. These numbers coincide with the published bounds of Eliahou and
Hercher; our contribution is the self-contained packaging and the identification
of the governing continued fraction with the one that drives the base-$6$
rotation underlying the dropping structure. We state clearly what is \emph{not}
proved: the nonexistence of cycles remains open, and our result is a lower bound,
conditional on the cited verification.

\medskip
\noindent\textbf{A standing invitation.} The author offers a prize of \textbf{US\$25}
to the first person who identifies a genuine mathematical error in the argument
of Sections \ref{sec:equation}--\ref{sec:bound} (see Section \ref{sec:prize}).
\end{abstract}

\section{Scope, and what is and is not claimed}

The Collatz map is $C(x) = x/2$ for even $x$ and $C(x) = 3x+1$ for odd $x$. A
\emph{nontrivial cycle} is a finite orbit of $C$ other than $\{1,2,4\}$. Whether
any nontrivial cycle exists is a famous open problem, and \textbf{nothing in this
note resolves it.} What we prove is a quantitative \emph{lower bound}: any
nontrivial cycle, should one exist, must be enormous. Such bounds are classical
(Crandall \cite{crandall}, Steiner \cite{steiner}, Eliahou \cite{eliahou},
Simons--de Weger \cite{simonsdeweger}, Hercher \cite{hercher}); our purpose is a
clean, elementary, and self-contained derivation in a framework that makes the
arithmetic transparent, and an honest delineation of the boundary between what
is proved and what is assumed.

\paragraph{Axioms and inputs.} We use only:
\begin{enumerate}[label=(A\arabic*),leftmargin=*]
\item \textbf{Elementary integer arithmetic} (the ring $\Z$, unique
factorisation, and the theory of continued fractions of a real number).
\item \textbf{The computational verification} that every integer $n < B := 2^{68}$
reaches $1$ under iteration of $C$ \cite{barina,oliveira}. Equivalently: the map
has no nontrivial cycle whose least element is below $B$. This is the only
non-elementary input, and our main theorem is \emph{conditional} on it.
\end{enumerate}
We do not invoke Baker's theorem on linear forms in logarithms; Section
\ref{sec:remarks} indicates how it removes the dependence on (A2) at the cost of
weaker constants.

\section{The odd map and the closed-word equation}\label{sec:equation}

It is convenient to track only the odd elements. For odd $n$, write
\[
  T(n) = \frac{3n+1}{2^{\,v(n)}}, \qquad v(n) = v_2(3n+1) \ge 1,
\]
where $v_2$ is the $2$-adic valuation; $T(n)$ is again odd. A cycle of $C$ through
odd values corresponds to a cycle of $T$, and conversely every cycle of $C$
(other than is forced through the even shadow) meets the odd numbers, so we lose
nothing by studying $T$.

\begin{lemma}[Closed-word equation]\label{lem:equation}
Let $n_0 \to n_1 \to \cdots \to n_{m-1} \to n_0$ be a cycle of $T$ of length
$m \ge 1$, with $v_i := v(n_i)$ and
\[
  K := \sum_{i=0}^{m-1} v_i, \qquad V_i := \sum_{j<i} v_j \ \ (V_0 = 0).
\]
Then
\begin{equation}\label{eq:cycle}
  \bigl(2^{K} - 3^{m}\bigr)\, n_0 \;=\; S, \qquad
  S \;=\; \sum_{i=0}^{m-1} 3^{\,m-1-i}\, 2^{\,V_i} \ \in \Z_{>0},
\end{equation}
and the same holds with $n_0$ replaced by any $n_j$ and the indices cyclically
shifted. In particular $2^{K} > 3^{m}$.
\end{lemma}

\begin{proof}
The defining relation is $2^{v_i} n_{i+1} = 3 n_i + 1$ (indices mod $m$).
Solving the linear recurrence around the loop, or equivalently substituting
repeatedly, gives $2^{K} n_0 = 3^{m} n_0 + S$ with $S$ as stated; this is the
standard cycle identity (see e.g.\ \cite{lagarias}, and \eqref{eq:product} below
for the multiplicative form). Each summand of $S$ is a positive integer, so
$S>0$, whence $2^K - 3^m = S/n_0 > 0$.
\end{proof}

\begin{remark}[The dictionary reading]
Equation \eqref{eq:cycle} is exactly the statement that the cycle is a
\emph{closed word} in the dropping dictionary. One round of $T$ on the residue
class fixing a parity word $w$ of $s$ odd steps and $k-s$ even steps is the exact
affine map $x \mapsto (3^{s}x + C_w)/2^{\,k-s}$, with $C_w$ the integer
``$+1$ accumulation'' \cite{affine}. Composing the words around a cycle multiplies
the slopes to $3^{m}/2^{K}$ and accumulates the intercepts to $S/2^{K}$; the fixed
point is $n_0 = S/(2^{K}-3^{m})$. The constant $S$ is the closed-word value of the
$C$'s. This viewpoint is not needed below, but it is the source of the framing.
\end{remark}

\section{The geometric--mean bound}\label{sec:bound}

The heart of the matter is a single inequality. Let $\nmin := \min_i n_i$.

\begin{lemma}[Least-element bound]\label{lem:gm}
For any cycle of $T$ of length $m$ with step sum $K$,
\begin{equation}\label{eq:gm}
  2^{K/m} \;\le\; 3 + \frac1{\nmin},
  \qquad\text{equivalently}\qquad
  \nmin \;\le\; \frac{1}{\,2^{K/m}-3\,}.
\end{equation}
Equality holds if and only if $m=1$ (the trivial cycle).
\end{lemma}

\begin{proof}
Dividing $2^{v_i} n_{i+1} = 3 n_i + 1$ by $n_i$ gives
$3 + \tfrac1{n_i} = 2^{v_i}\, n_{i+1}/n_i$. Taking the product over a full loop,
the factors $n_{i+1}/n_i$ telescope to $1$, so
\begin{equation}\label{eq:product}
  \prod_{i=0}^{m-1}\Bigl(3 + \frac1{n_i}\Bigr) \;=\; 2^{\,K}.
\end{equation}
Since $n_i \ge \nmin$ for every $i$, each factor satisfies
$3 + 1/n_i \le 3 + 1/\nmin$, hence
$2^{K} \le (3 + 1/\nmin)^{m}$. Taking $m$-th roots gives \eqref{eq:gm}.
Because $2^{K}>3^{m}$ (Lemma \ref{lem:equation}), we have $2^{K/m}>3$, so the
right-hand form is well defined and positive. Equality in the product step
requires all $n_i$ equal, i.e.\ $m=1$, which is the cycle $\{1\}$
($3\cdot1+1 = 2^{2}$, so $K=2$, and \eqref{eq:gm} reads $2^{2}= 4 = 3+1$).
\end{proof}

\noindent
The inequality \eqref{eq:gm} is exact and elementary; it is the only analytic
ingredient. Everything else is arithmetic of $\log_2 3$.

\section{The forced approximation and the length bound}\label{sec:main}

Write $\theta := \log_2 3 = 1.5849625\ldots$, an irrational number. Combining the
least-element bound with the verification (A2) pins the step ratio $K/m$ into a
microscopic interval just above $\theta$.

\begin{lemma}[Forced approximation]\label{lem:forced}
Assume (A2). For any nontrivial cycle of $T$,
\begin{equation}\label{eq:forced}
  \theta \;<\; \frac{K}{m} \;\le\; \log_2\!\Bigl(3 + \tfrac1{B}\Bigr)
  \;=\; \theta + \delta, \qquad
  0 < \delta \le \frac{1}{B\,\ln 8} < 2^{-69}.
\end{equation}
\end{lemma}

\begin{proof}
A nontrivial cycle is not $\{1,2,4\}$, so by (A2) its least element satisfies
$\nmin \ge B$. Lemma \ref{lem:gm} gives $2^{K/m} \le 3 + 1/\nmin \le 3 + 1/B$,
i.e.\ $K/m \le \log_2(3+1/B)$. The lower bound $K/m>\theta$ is $2^{K}>3^{m}$.
Finally
\[
  \delta = \log_2\!\Bigl(\tfrac{3+1/B}{3}\Bigr)
         = \log_2\!\Bigl(1 + \tfrac1{3B}\Bigr)
         \le \frac{1}{3B\,\ln 2} = \frac1{B\,\ln 8},
\]
using $\ln(1+x)\le x$; numerically $1/(B\ln 8) = 2^{-69.06\ldots}$.
\end{proof}

The remaining question is purely Diophantine: \emph{what is the smallest
denominator $m$ of a rational $K/m$ lying in the interval
$(\theta,\ \theta+\delta]$?} This is answered exactly by the continued-fraction
(Stern--Brocot) algorithm for the simplest fraction in an interval.

\begin{lemma}[Smallest denominator in an interval]\label{lem:simplest}
Let $0<\delta$ and let $p/q$ be the unique fraction of least denominator lying in
$(\theta,\theta+\delta]$. Then, by definition of least denominator, every rational
in that interval has denominator at least $q$. For $\delta = 2^{-69.06}$
(the value of Lemma \ref{lem:forced} at $B=2^{68}$),
\[
  \frac{p}{q} \;=\; \frac{114{,}208{,}327{,}604}{72{,}057{,}431{,}991},
  \qquad q = 72{,}057{,}431{,}991 \approx 2^{36.07}.
\]
\end{lemma}

\begin{proof}
The fraction of least denominator in a real interval is computed by the classical
mediant/continued-fraction descent and is unique; any other fraction in the
interval has strictly larger denominator (see \cite[Ch.~I]{khinchin} or
\cite{stern-brocot}). The displayed value is the output of that algorithm applied
to the endpoints $\theta$ and $\theta+\delta$, evaluated to $80$ decimal digits;
both endpoints and the membership $\theta < p/q \le \theta+\delta$ are verified at
that precision. The numerator/denominator are consecutive in the continued
fraction of $\theta$, whose partial quotients begin
$[1;1,1,2,2,3,1,5,2,23,2,2,1,1,55,\ldots]$.
\end{proof}

\begin{theorem}[Cycle lower bound]\label{thm:main}
Assume (A2). Every nontrivial cycle of the $3x+1$ map has
\[
  m \ \ge\ 72{,}057{,}431{,}991 \ \text{(odd elements)}, \qquad
  K \ \ge\ 114{,}208{,}327{,}604 \ \text{(halvings)},
\]
hence a full period
\[
  \lambda = m + K \ \ge\ 186{,}265{,}759{,}595 \ \approx\ 1.86\times10^{11}
  \ \text{steps},
\]
and least element $\nmin \ge 2^{68} \approx 2.95\times10^{20}$.
\end{theorem}

\begin{proof}
By Lemma \ref{lem:forced}, $K/m \in (\theta,\theta+\delta]$ with
$\delta = 2^{-69.06}$. Writing $K/m$ in lowest terms $p'/q'$ (so $q' \mid m$),
Lemma \ref{lem:simplest} gives $q' \ge q = 72{,}057{,}431{,}991$, whence
$m \ge q' \ge q$. Then $K > m\theta \ge q\,\theta$ gives the stated bound on $K$
(the exact value is the matching numerator), and $\lambda = m+K$. The bound
$\nmin \ge B = 2^{68}$ is (A2). The full period counts one $3x+1$ step and $v_i$
halvings per odd element, i.e.\ $\sum_i (1+v_i) = m+K$.
\end{proof}

\begin{remark}
The numerator $K = 114{,}208{,}327{,}604$ coincides with the cycle--length lower
bound reported by Hercher \cite{hercher} under the same verification height, and
Eliahou's earlier bound \cite{eliahou} arises identically from a smaller $B$.
This agreement is a strong external check: our elementary route reproduces the
state of the art exactly.
\end{remark}

\section{The continued fraction is the rotation's}\label{sec:cf}

The denominator $q$ in Theorem \ref{thm:main} is a convergent denominator of
$\theta = \log_2 3$, whose convergents are
\[
  \tfrac11,\ \tfrac21,\ \tfrac32,\ \tfrac85,\ \tfrac{19}{12},\ \tfrac{65}{41},\
  \tfrac{84}{53},\ \tfrac{485}{306},\ \tfrac{1054}{665},\ \ldots
\]
This is the same continued fraction that governs the base-$6$ logarithmic
rotation $\theta \mapsto \theta + \log_6 3$ under which every Collatz step is a
rigid turn: with $\alpha = \log_6 3 = \theta/(1+\theta)$, the two share a
continued-fraction tail, and the convergent denominators of $\alpha$ are
$1,1,2,3,5,13,31,106,137,\ldots$\,---\,the near-return periods of the rotation,
the $13/31/137$ spiral-arm counts of an orbit plot, and the resonant levels of
the dropping alphabet \cite{droppingdict}. The arithmetic that makes cycles rare
(a good rational approximation of $\theta$ is expensive) is the same arithmetic
that makes the rotation quasiperiodic. The cycle bound and the wobble's $44$-step
quasi-period are two consequences of one irrational.

\section{What this does not prove, and how (A2) could be removed}\label{sec:remarks}

\paragraph{Not a nonexistence proof.} Theorem \ref{thm:main} bounds how large a
cycle must be; it does not exclude cycles. The obstruction to ``no cycles'' is
that \eqref{eq:gm} controls only the least element through the step ratio $K/m$;
a hypothetical cycle with $K/m$ an extraordinarily good rational approximation to
$\theta$ (denominator beyond the verification's reach) is not contradicted by any
argument here. Closing that gap is exactly the open problem.

\paragraph{Removing the verification.} The dependence on (A2) is removable in
principle. Baker's theorem on linear forms in logarithms gives an effective lower
bound $|K\ln 2 - m\ln 3| \ge K^{-\kappa}\,c$ for explicit $c,\kappa$, which bounds
$\delta$ from below in terms of $m$ alone and yields an \emph{unconditional}
(if numerically weaker, and structurally more involved) lower bound on cycle
length; this is the route of Simons--de Weger \cite{simonsdeweger}. We have kept
(A2) because it is elementary to state, currently verified well beyond
$2^{68}$, and gives the sharper constant.

\paragraph{Axiom audit.} The chain uses (i) integer arithmetic and continued
fractions, used in Lemmas \ref{lem:equation}--\ref{lem:simplest} and certain at
the stated $80$-digit precision; and (ii) the verification (A2), an exhaustive
but finite computation whose correctness is the one empirical assumption. We are
confident in (i) unconditionally and in the theorem conditional on (ii).

\section{An invitation to find a flaw}\label{sec:prize}

We are sufficiently confident in the derivation of
Sections \ref{sec:equation}--\ref{sec:bound} (Lemmas
\ref{lem:equation}--\ref{lem:simplest} and Theorem \ref{thm:main}) to back it:
\textbf{the author will pay US\$25 to the first person who exhibits a genuine
mathematical error} in those statements or their proofs\,---\,an incorrect
inequality, a gap in the continued-fraction argument, a miscomputed constant, or
a misuse of the verification (A2). Cosmetic, expository, or
``the result is already known'' objections do not qualify (we cite the prior art
in Section \ref{sec:main}); the claim on offer is correctness, not novelty.
Reports may be sent to the author.

\begin{thebibliography}{9}

\bibitem{barina} D.\ Barina, \emph{Convergence verification of the Collatz
problem}, J.\ Supercomput.\ \textbf{77} (2021), 2681--2688.

\bibitem{oliveira} T.\ Oliveira e Silva, \emph{Empirical verification of the
$3x+1$ and related conjectures}, in \emph{The Ultimate Challenge: The $3x+1$
Problem}, AMS, 2010.

\bibitem{crandall} R.\ E.\ Crandall, \emph{On the $3x+1$ problem}, Math.\ Comp.\
\textbf{32} (1978), 1281--1292.

\bibitem{steiner} R.\ P.\ Steiner, \emph{A theorem on the Syracuse problem},
Proc.\ 7th Manitoba Conf.\ Numer.\ Math., 1977, 553--559.

\bibitem{eliahou} S.\ Eliahou, \emph{The $3x+1$ problem: new lower bounds on
nontrivial cycle lengths}, Discrete Math.\ \textbf{118} (1993), 45--56.

\bibitem{simonsdeweger} J.\ Simons and B.\ de Weger, \emph{Theoretical and
computational bounds for $m$-cycles of the $3n+1$ problem}, Acta Arith.\
\textbf{117} (2005), 51--70.

\bibitem{hercher} C.\ Hercher, \emph{There are no Collatz cycles with fewer than
$\sim 1.14\times 10^{11}$ terms} (cycle-length lower bound under current
verification), 2023.

\bibitem{lagarias} J.\ C.\ Lagarias, \emph{The $3x+1$ problem and its
generalizations}, Amer.\ Math.\ Monthly \textbf{92} (1985), 3--23.

\bibitem{khinchin} A.\ Ya.\ Khinchin, \emph{Continued Fractions}, Univ.\ of
Chicago Press, 1964.

\bibitem{stern-brocot} R.\ L.\ Graham, D.\ E.\ Knuth, O.\ Patashnik,
\emph{Concrete Mathematics}, 2nd ed., Addison--Wesley, 1994, \S4.5.

\bibitem{affine} \emph{Affine Orbit Structure} and \emph{Nested Dropping Sets},
research notes, Collatz exploration repository, 2026.

\bibitem{droppingdict} \emph{The Dropping Dictionary}, whycollatzworks.com /
Collatz exploration repository, 2026.

\end{thebibliography}

\end{document}
